\section{Resultados}

% Las cifras de este capitulo se inyectan desde los CSV del estudio de
% ablacion, mediante las tablas generadas en docs/tablas/. No escribir valores
% numericos a mano: regenerar con
%     python -m src.experiments.ablation
% y volver a compilar.

\subsection{Comparacion entre representaciones}

\begin{table}[htbp]
\caption{Exactitud por representacion bajo validacion cruzada}
\label{tab:resultados-via}
\begin{tabular}{llrrr}
\toprule
Via & Representacion & $M$ exactitud & $DE$ & $M$ macro-F1 \\
\midrule
A & Fisicas & 0.841 & 0.015 & 0.841 \\
B & Eigen & 0.900 & 0.011 & 0.899 \\
\bottomrule
\end{tabular}
\note{Validacion cruzada estratificada de 5 pliegues sobre el mismo
clasificador de minimos cuadrados con $\lambda = 0.01$ y $k = 50$.
Elaboracion propia.}
\end{table}


% TODO: describir la tabla en prosa.

\begin{figure}[htbp]
\caption{Exactitud por representacion bajo validacion cruzada}
\figurascoal{ablacion_via.png}{0.6\textwidth}
\note{Media y desviacion estandar sobre cinco pliegues estratificados.
Elaboracion propia.}
\end{figure}

\subsection{Efecto del numero de componentes}

\begin{figure}[htbp]
\caption{Exactitud y varianza explicada en funcion de $k$}
\figurascoal{ablacion_k.png}{\textwidth}
\note{Vias B y C. La linea discontinua marca el noventa por ciento de
varianza explicada acumulada. Elaboracion propia.}
\end{figure}

% TODO: comentar el punto de saturacion.

\subsection{Efecto de la regularizacion}

\begin{figure}[htbp]
\caption{Exactitud y numero de condicion en funcion de $\lambda$}
\figurascoal{ablacion_lambda.png}{\textwidth}
\note{El panel derecho muestra $\operatorname{cond}_2(X^{T}X+\lambda I)$ en
escala logaritmica. Elaboracion propia.}
\end{figure}

% TODO: comentar el optimo y la caida del numero de condicion.

\subsection{Comparacion contra los baselines}

\begin{table}[htbp]
\caption{Comparacion contra los baselines}
\label{tab:baselines}
\begin{tabular}{llr}
\toprule
Baseline & Via & Exactitud \\
\midrule
aleatorio & - & 0.100 \\
knn\_k5 & A & 0.927 \\
minimos\_cuadrados & A & 0.821 \\
knn\_k5 & B & 0.976 \\
minimos\_cuadrados & B & 0.872 \\
\bottomrule
\end{tabular}
\note{Todos los valores sobre la misma particion de prueba.
Elaboracion propia.}
\end{table}


% TODO: comentar la brecha entre la via C y la cabeza original de MobileNetV3.

\subsection{Analisis cualitativo}

\begin{figure}[htbp]
\caption{Primeros dieciseis eigen-objetos}
\figurascoal{eigenobjetos.png}{0.6\textwidth}
\note{Cada panel indica la fraccion de varianza explicada. Elaboracion propia.}
\end{figure}

\begin{figure}[htbp]
\caption{Proyeccion sobre las dos primeras componentes principales}
\figurascoal{scatter_pca.png}{0.6\textwidth}
\note{Color por clase. Elaboracion propia.}
\end{figure}

\begin{figure}[htbp]
\caption{Matriz de confusion de la via A}
\figurascoal{confusion_A.png}{0.55\textwidth}
\note{Filas normalizadas por clase real. Elaboracion propia.}
\end{figure}

\begin{figure}[htbp]
\caption{Matriz de confusion de la via B}
\figurascoal{confusion_B.png}{0.55\textwidth}
\note{Filas normalizadas por clase real. Elaboracion propia.}
\end{figure}

\begin{figure}[htbp]
\caption{Matriz de confusion de la via C}
\figurascoal{confusion_C.png}{0.55\textwidth}
\note{Filas normalizadas por clase real. Elaboracion propia.}
\end{figure}
